% libre-compendium
%
% Copyright (C) 2026 Richard Qin
%
% This document is dual-licensed under either:
% - GNU Free Documentation License, Version 1.3 or any later version
%   published by the Free Software Foundation; with no Invariant Sections,
%   no Front-Cover Texts, and no Back-Cover Texts (see COPYING.DOC); or
% - Creative Commons Attribution-ShareAlike 4.0 International License
%   (see LICENSE-CC-BY-SA-4.0).

\chapter{矩阵微分、广义逆与正交投影}

本章默认所有变量为实变量。向量均写成列向量，标量函数 $f:\R^n\to\R$ 的梯度为列向量，满足
\[
df=(\nabla f(x))^{\mathsf T}dx.
\]
对矩阵变量 $X\in\R^{m\times n}$，矩阵梯度由
\[
df=\tr\!\left((\nabla_Xf)^{\mathsf T}dX\right)
=\langle\nabla_Xf,dX\rangle_F
\]
定义。这两个等式是本章公式方向和维数的最终判定标准。

\section{从全微分到梯度与 Jacobian}

\begin{definition}[方向导数与 Fr\'echet 导数]
设 $F:\R^n\to\R^m$。若存在线性映射 $L:\R^n\to\R^m$ 使
\[
F(x+h)=F(x)+Lh+o(\lVert h\rVert),
\qquad h\to0,
\]
则称 $F$ 在 $x$ 处可微，$L$ 称为 $F$ 在 $x$ 处的导数。在标准基下，$L$ 的矩阵是 Jacobian
\[
J_F(x)=\left(\frac{\partial F_i}{\partial x_j}\right)_{m\times n},
\qquad dF=J_F(x)\,dx.
\]
\end{definition}

文献对“分子布局”和“分母布局”的命名并不完全统一。因此，不应只凭布局名称判断公式，而应同时写出矩阵尺寸和微分关系。本章始终使用 $J_F\in\R^{m\times n}$ 且 $dF=J_Fdx$。

\begin{proposition}[链式法则]
设 $F:\R^n\to\R^m$、$G:\R^m\to\R^p$ 可微，则
\[
J_{G\circ F}(x)=J_G(F(x))J_F(x).
\]
特别地，若 $f:\R^m\to\R$ 是标量函数，则
\[
\nabla_x(f\circ F)=J_F(x)^{\mathsf T}\nabla_yf\big|_{y=F(x)}.
\]
\end{proposition}
\begin{proof}
由可微性，
\[
dF=J_Fdx,
\qquad
dG=J_GdF.
\]
代入得 $d(G\circ F)=J_GJ_Fdx$。对标量函数，
\[
df=(\nabla_yf)^{\mathsf T}dy
=(\nabla_yf)^{\mathsf T}J_Fdx
=(J_F^{\mathsf T}\nabla_yf)^{\mathsf T}dx,
\]
比较 $dx$ 的系数即得梯度公式。
\end{proof}

\begin{example}[title={仿射变换后的逐元素函数}]
设 $a=Wx+b$，$z=\sigma(a)$，其中 $\sigma$ 逐元素作用。求 $z$ 对 $x$ 的 Jacobian。
\begin{solution}
有
\[
da=Wdx,
\qquad
dz=D_\sigma(a)da,
\]
其中
\[
D_\sigma(a)=\diag\bigl(\sigma'(a_1),\ldots,\sigma'(a_m)\bigr).
\]
因此
\[
J_z(x)=D_\sigma(a)W.
\]
该式的尺寸为 $m\times n$，与 $z\in\R^m$、$x\in\R^n$ 一致。
\end{solution}
\end{example}

\section{矩阵微分与迹技巧}

\begin{definition}[矩阵变量的梯度]
设 $f:\R^{m\times n}\to\R$。若
\[
df=\sum_{i,j}\frac{\partial f}{\partial x_{ij}}dx_{ij},
\]
则定义
\[
\nabla_Xf=\left(\frac{\partial f}{\partial x_{ij}}\right)_{m\times n}.
\]
由 Frobenius 内积的定义，
\[
df=\tr\!\left((\nabla_Xf)^{\mathsf T}dX\right).
\]
\end{definition}

矩阵求导的稳定流程是：先求微分，再利用迹的循环性把 $dX$ 移到末尾，最后与
$df=\tr(G^{\mathsf T}dX)$ 比较，得到 $G=\nabla_Xf$。

\begin{proposition}[基本微分法则]
在矩阵乘法尺寸相容时，
\[
d(X+Y)=dX+dY,
\qquad
d(XY)=(dX)Y+X(dY),
\]
\[
d(X^{\mathsf T})=(dX)^{\mathsf T},
\qquad
d\tr(X)=\tr(dX).
\]
若 $X$ 可逆，则
\[
d(X^{-1})=-X^{-1}(dX)X^{-1}.
\]
\end{proposition}
\begin{proof}
乘法公式来自
\[
(X+dX)(Y+dY)-XY=(dX)Y+X(dY)+(dX)(dY),
\]
其中 $(dX)(dY)$ 是二阶小量。对逆矩阵，对 $XX^{-1}=I$ 求微分：
\[
(dX)X^{-1}+X\,d(X^{-1})=0.
\]
左乘 $X^{-1}$ 即得结论。
\end{proof}

\begin{proposition}[迹的常用恒等式]
若尺寸相容，则
\[
\tr(AB)=\tr(BA),
\qquad
\tr(ABC)=\tr(BCA)=\tr(CAB),
\]
\[
\tr(A^{\mathsf T}B)=\langle A,B\rangle_F,
\qquad
\tr(A\,dX^{\mathsf T})=\tr(A^{\mathsf T}dX).
\]
迹只能循环换位，不能任意交换两个相邻因子。
\end{proposition}

\begin{theorem}[常用矩阵梯度]
设 $A,B,C$ 为常矩阵，则
\[
\nabla_X\tr(A^{\mathsf T}X)=A,
\]
\[
\nabla_X\tr(AXB)=A^{\mathsf T}B^{\mathsf T},
\]
\[
\nabla_X\tr(X^{\mathsf T}AX)=(A+A^{\mathsf T})X,
\]
而双侧二次型的完整公式是
\[
\boxed{\nabla_X\tr(X^{\mathsf T}AXB)=AXB^{\mathsf T}+A^{\mathsf T}XB.}
\]
特别地，若 $A,B$ 均对称，则梯度为 $2AXB$。
\end{theorem}
\begin{proof}
以最后一式为例。令 $f=\tr(X^{\mathsf T}AXB)$，则
\begin{align*}
df
&=\tr((dX)^{\mathsf T}AXB)+\tr(X^{\mathsf T}A(dX)B)\\
&=\tr((AXB^{\mathsf T})^{\mathsf T}dX)
 +\tr((A^{\mathsf T}XB)^{\mathsf T}dX).
\end{align*}
比较系数得
\[
\nabla_Xf=AXB^{\mathsf T}+A^{\mathsf T}XB.
\]
其余公式均可用同一方法推导。
\end{proof}

\begin{pitfall}[不要直接套用未经验证的公式]
矩阵乘法不可交换，对称性也不能默认。例如
$\nabla_x(x^{\mathsf T}Ax)=(A+A^{\mathsf T})x$，只有 $A=A^{\mathsf T}$ 时才等于 $2Ax$。检查公式时至少应核对尺寸，并通过微分恒等式复核转置位置。
\end{pitfall}

\section{二次型、Frobenius 范数与最小二乘}

\begin{proposition}[向量二次型]
对 $f(x)=x^{\mathsf T}Ax+b^{\mathsf T}x+c$，
\[
\nabla f(x)=(A+A^{\mathsf T})x+b,
\qquad
\nabla^2f(x)=A+A^{\mathsf T}.
\]
若 $A$ 对称，则 $f$ 严格凸当且仅当 $A$ 正定。
\end{proposition}

\begin{theorem}[Frobenius 范数型目标]
设
\[
f(X)=\frac12\lVert AXB-C\rVert_F^2.
\]
则
\[
\nabla_Xf=A^{\mathsf T}(AXB-C)B^{\mathsf T}.
\]
\end{theorem}
\begin{proof}
令 $R=AXB-C$，则 $dR=A(dX)B$。由
$f=\frac12\tr(R^{\mathsf T}R)$，
\begin{align*}
df
&=\tr(R^{\mathsf T}dR)
=\tr(R^{\mathsf T}A(dX)B)\\
&=\tr\!\left((A^{\mathsf T}RB^{\mathsf T})^{\mathsf T}dX\right).
\end{align*}
因此梯度如上。
\end{proof}

\begin{corollary}[线性最小二乘的正规方程]
对
\[
f(x)=\frac12\lVert Ax-b\rVert_2^2,
\]
有
\[
\nabla f(x)=A^{\mathsf T}(Ax-b),
\qquad
\nabla^2f=A^{\mathsf T}A.
\]
因此任一最小二乘解都满足正规方程
\[
A^{\mathsf T}Ax=A^{\mathsf T}b.
\]
若 $A$ 列满秩，则 $A^{\mathsf T}A$ 正定，最小二乘解唯一：
\[
x_*=(A^{\mathsf T}A)^{-1}A^{\mathsf T}b.
\]
\end{corollary}

\begin{remark}
正规方程在理论推导中自然，但数值计算中会把条件数平方：
$\kappa_2(A^{\mathsf T}A)=\kappa_2(A)^2$。实际求解常优先使用 QR 分解或 SVD。
\end{remark}

\section{加权与约束最小二乘}

\begin{definition}[加权最小二乘]
设 $W\in\R^{m\times m}$ 对称正定。加权最小二乘问题为
\[
\min_x\frac12(Ax-b)^{\mathsf T}W(Ax-b)
=\min_x\frac12\lVert Ax-b\rVert_W^2,
\]
其中 $\langle u,v\rangle_W=u^{\mathsf T}Wv$ 是 $\R^m$ 上的内积。
\end{definition}

\begin{theorem}[加权正规方程]
加权最小二乘解满足
\[
A^{\mathsf T}WAx=A^{\mathsf T}Wb.
\]
等价地，残差 $r=b-Ax$ 在 $W$-内积下与 $\operatorname{Im}A$ 正交：
\[
A^{\mathsf T}Wr=0.
\]
若 $A$ 列满秩，则解唯一，且
\[
x_*=(A^{\mathsf T}WA)^{-1}A^{\mathsf T}Wb.
\]
\end{theorem}
\begin{proof}
令 $f(x)=\frac12(Ax-b)^{\mathsf T}W(Ax-b)$。由 $W=W^{\mathsf T}$，
\[
df=(Ax-b)^{\mathsf T}WA,dx,
\]
故 $\nabla f=A^{\mathsf T}W(Ax-b)$。驻点条件给出加权正规方程。若 $A$ 列满秩，则对 $h\ne0$，
\[
h^{\mathsf T}A^{\mathsf T}WAh=(Ah)^{\mathsf T}W(Ah)>0,
\]
所以 $A^{\mathsf T}WA$ 正定可逆。
\end{proof}

\begin{corollary}[加权正交投影]
若 $A$ 列满秩，则
\[
P_W=A(A^{\mathsf T}WA)^{-1}A^{\mathsf T}W
\]
满足 $P_W^2=P_W$，并把 $b$ 投影到 $\operatorname{Im}A$。它一般不满足 $P_W^{\mathsf T}=P_W$，但满足
\[
P_W^{\mathsf T}W=WP_W,
\]
因而它是 $W$-内积下的正交投影。
\end{corollary}

\begin{definition}[等式约束最小二乘]
考虑
\[
\min_x\frac12\lVert Ax-b\rVert_2^2,
\qquad Cx=d.
\]
其 Lagrange 函数为
\[
\mathcal L(x,\lambda)
=\frac12\lVert Ax-b\rVert_2^2
+\lambda^{\mathsf T}(Cx-d).
\]
\end{definition}

\begin{theorem}[等式约束的 KKT 方程]
任一约束最小二乘解都与某个乘子 $\lambda$ 一起满足
\[
\begin{pmatrix}
A^{\mathsf T}A&C^{\mathsf T}\\
C&0
\end{pmatrix}
\begin{pmatrix}x\\\lambda\end{pmatrix}
=
\begin{pmatrix}A^{\mathsf T}b\\d\end{pmatrix}.
\]
若 $C$ 行满秩且
\[
\ker A\cap\ker C=\{0\},
\]
则 KKT 矩阵可逆，约束最小二乘解唯一。
\end{theorem}
\begin{proof}
对 $x$ 和 $\lambda$ 分别求梯度，得
\[
A^{\mathsf T}(Ax-b)+C^{\mathsf T}\lambda=0,
\qquad Cx=d,
\]
这就是所给分块方程。

为证可逆性，设齐次 KKT 方程有解 $(h,\mu)$。左乘第一个方程以 $h^{\mathsf T}$，并利用 $Ch=0$，得
\[
\lVert Ah\rVert_2^2=0.
\]
因此 $h\in\ker A\cap\ker C$，从而 $h=0$。随后 $C^{\mathsf T}\mu=0$，由 $C$ 行满秩得 $\mu=0$。齐次核为零，故 KKT 矩阵可逆。
\end{proof}

\section{行列式、逆矩阵与对数行列式}

\begin{theorem}[Jacobi 行列式公式]
若 $X$ 可逆，则
\[
d\det X=\det(X)\tr(X^{-1}dX),
\]
从而
\[
\nabla_X\det X=\det(X)X^{-\mathsf T}.
\]
若 $X$ 正定，则
\[
d\log\det X=\tr(X^{-1}dX),
\qquad
\nabla_X\log\det X=X^{-\mathsf T}=X^{-1}.
\]
\end{theorem}
\begin{proof}
对充分小的 $dX$，
\[
\det(X+dX)=\det X\,\det(I+X^{-1}dX).
\]
而 $\det(I+E)=1+\tr(E)+o(\lVert E\rVert)$，因此
\[
\det(X+dX)-\det X
=\det X\,\tr(X^{-1}dX)+o(\lVert dX\rVert).
\]
再用
$\tr(X^{-1}dX)=\tr((X^{-\mathsf T})^{\mathsf T}dX)$
读出梯度。
\end{proof}

\begin{remark}
$\log\det X$ 在实数范围内通常只对正定矩阵使用。对一般可逆实矩阵，可考虑 $\log|\det X|$，其微分在行列式不过零的连通区域内仍为 $\tr(X^{-1}dX)$。
\end{remark}

\begin{example}[title={带对数行列式的梯度}]
设 $X$ 正定，
\[
f(X)=\tr(CX^{-1})-\alpha\log\det X.
\]
求 $\nabla_Xf$。
\begin{solution}
由逆矩阵微分公式，
\begin{align*}
d\tr(CX^{-1})
&=-\tr(CX^{-1}(dX)X^{-1})\\
&=-\tr(X^{-1}CX^{-1}dX).
\end{align*}
因此
\[
\nabla_Xf
=-X^{-\mathsf T}C^{\mathsf T}X^{-\mathsf T}
-\alpha X^{-\mathsf T}.
\]
若 $X,C$ 均对称，则化为
$-X^{-1}CX^{-1}-\alpha X^{-1}$。
\end{solution}
\end{example}

\section{对称矩阵特征值的微分}

\begin{theorem}[简单特征值的一阶扰动]
设 $A=A^{\mathsf T}$，$\lambda_i$ 是 $A$ 的单重特征值，$u_i$ 是对应的单位特征向量。对对称扰动 $dA$，
\[
d\lambda_i=u_i^{\mathsf T}(dA)u_i
=\tr(u_iu_i^{\mathsf T}dA).
\]
因此在对称矩阵空间中
\[
\nabla_A\lambda_i=u_iu_i^{\mathsf T}.
\]
\end{theorem}
\begin{proof}
对特征方程 $Au_i=\lambda_iu_i$ 求微分：
\[
(dA)u_i+A,du_i
=(d\lambda_i)u_i+\lambda_i,du_i.
\]
左乘 $u_i^{\mathsf T}$。由 $A$ 对称，
$u_i^{\mathsf T}A=\lambda_i u_i^{\mathsf T}$，两边含 $du_i$ 的项相消；再用 $u_i^{\mathsf T}u_i=1$，得
\[
d\lambda_i=u_i^{\mathsf T}(dA)u_i.
\]
\end{proof}

\begin{pitfall}[重特征值处不能指定唯一梯度]
当 $\lambda$ 是重特征值时，特征子空间内的单位特征向量不唯一，$uu^{\mathsf T}$ 也随选择改变。例如最大特征值函数在顶部特征值重复时一般不可微；此时应使用特征子空间上的扰动分析或次梯度。
\end{pitfall}

\begin{corollary}[最大特征值的梯度]
若 $A=A^{\mathsf T}$ 的最大特征值 $\lambda_{\max}(A)$ 为单重，对应单位特征向量为 $u$，则
\[
\nabla_A\lambda_{\max}(A)=uu^{\mathsf T}.
\]
\end{corollary}

\section{Vec 算子、Kronecker 积与矩阵 Jacobian}

\begin{definition}[Vec 与 Kronecker 积]
对 $X=(x_1,\ldots,x_n)\in\R^{m\times n}$，定义按列展开
\[
\operatorname{vec}(X)
=(x_1^{\mathsf T},\ldots,x_n^{\mathsf T})^{\mathsf T}\in\R^{mn}.
\]
对 $A=(a_{ij})\in\R^{p\times q}$、$B\in\R^{r\times s}$，定义
\[
A\otimes B=(a_{ij}B)\in\R^{pr\times qs}.
\]
\end{definition}

\begin{theorem}[Vec 恒等式]
若尺寸相容，则
\[
\boxed{\operatorname{vec}(AXB)
=(B^{\mathsf T}\otimes A)\operatorname{vec}(X).}
\]
因此对 $Y=AXB+C$，
\[
\frac{\partial\operatorname{vec}(Y)}
{\partial\operatorname{vec}(X)^{\mathsf T}}
=B^{\mathsf T}\otimes A.
\]
\end{theorem}
\begin{proof}
记 $B=(b_1,\ldots,b_s)$。$AXB$ 的第 $j$ 列为
\[
AXb_j=\sum_{k=1}^q b_{kj}Ax_k.
\]
把所有列依次堆叠，每一个块对 $x_k$ 的系数正是 $b_{kj}A$，这与 $B^{\mathsf T}\otimes A$ 的块结构一致。
\end{proof}

\begin{corollary}[Sylvester 方程的向量化]
矩阵方程
\[
AX+XB=C
\]
等价于
\[
\bigl(I\otimes A+B^{\mathsf T}\otimes I\bigr)
\operatorname{vec}(X)=\operatorname{vec}(C).
\]
因此，若 $A$ 的任一特征值与 $-B$ 的任一特征值都不相同，则 Sylvester 方程有唯一解。
\end{corollary}

\section{对称矩阵变量的约束微分}

若 $X$ 被限制为对称矩阵，则允许的扰动必须满足 $dX=dX^{\mathsf T}$。假设先把 $X$ 当作一般矩阵求得
\[
df=\tr(G^{\mathsf T}dX),
\]
则对称约束下
\[
df=\tr\!\left(\left(\frac{G+G^{\mathsf T}}2\right)^{\mathsf T}dX\right).
\]
因此只有 $G$ 的对称部分会作用在可行方向上。

\begin{pitfall}[对称约束与独立参数求导不是同一种记号]
把 $X$ 视为对称矩阵时，$x_{ij}$ 与 $x_{ji}$ 不是两个独立变量。若改用上三角元素组成的参数向量，非对角参数的系数会合并两个对称位置的贡献。因此必须先说明是“Frobenius 内积下的可行梯度”，还是“对独立上三角参数的偏导数”。
\end{pitfall}

\section{一般广义逆}

\begin{definition}[一般广义逆]
对 $A\in\R^{m\times n}$，若 $G\in\R^{n\times m}$ 满足
\[
AGA=A,
\]
则称 $G$ 是 $A$ 的一个广义逆，常记为 $A^{-}$。一般广义逆通常不唯一。
\end{definition}

\begin{theorem}[广义逆的存在性]
每个实矩阵都存在广义逆。
\end{theorem}
\begin{proof}
设 $\rank A=r$。由矩阵等价标准形，存在可逆矩阵 $P,Q$ 使
\[
A=PE_{m,n}Q,
\qquad
E_{m,n}=\begin{pmatrix}I_r&0\\0&0\end{pmatrix}\in\R^{m\times n}.
\]
令
\[
G=Q^{-1}E_{n,m}P^{-1},
\qquad E_{n,m}=E_{m,n}^{\mathsf T},
\]
则直接相乘得 $AGA=A$。
\end{proof}

\begin{theorem}[相容线性方程的通解]
若 $Ax=b$ 有解，$A^{-}$ 是 $A$ 的任意广义逆，则
\[
x=A^{-}b+(I-A^{-}A)z,
\qquad z\in\R^n,
\]
给出 $Ax=b$ 的全部解。
\end{theorem}
\begin{proof}
由 $b\in\operatorname{Im}A$，可写成 $b=Au$。因此
\[
AA^{-}b=AA^{-}Au=Au=b.
\]
对任意 $z$，
\[
A\bigl(A^{-}b+(I-A^{-}A)z\bigr)
=b+(A-AA^{-}A)z=b.
\]
反之，若 $Ax=b$，则
\[
x=A^{-}Ax+(I-A^{-}A)x
=A^{-}b+(I-A^{-}A)x,
\]
所以取 $z=x$ 即可包含任意解。
\end{proof}

\begin{remark}
条件“$Ax=b$ 有解”不能删去。对不相容方程，$A^{-}b$ 未必是最小二乘解，甚至 $AA^{-}b$ 也未必等于 $b$。要得到唯一且具有正交几何意义的逆，需要 Moore--Penrose 条件。
\end{remark}

\section{Moore--Penrose 逆与 SVD 构造}

\begin{definition}[Moore--Penrose 逆]
矩阵 $A\in\R^{m\times n}$ 的 Moore--Penrose 逆是满足
\[
AA^+A=A,
\qquad
A^+AA^+=A^+,
\]
\[
(AA^+)^{\mathsf T}=AA^+,
\qquad
(A^+A)^{\mathsf T}=A^+A
\]
的矩阵 $A^+\in\R^{n\times m}$。
\end{definition}

\begin{theorem}[Moore--Penrose 逆的存在唯一性与 SVD 公式]
若
\[
A=U\Sigma V^{\mathsf T}
\]
是 $A$ 的 SVD，且非零奇异值为
$\sigma_1,\ldots,\sigma_r$，则
\[
A^+=V\Sigma^+U^{\mathsf T},
\]
其中 $\Sigma^+$ 把每个非零奇异值换成倒数，并转置矩形尺寸。该矩阵是唯一的。
\end{theorem}
\begin{proof}
对角块上有
\[
\Sigma\Sigma^+\Sigma=\Sigma,
\qquad
\Sigma^+\Sigma\Sigma^+=\Sigma^+,
\]
并且 $\Sigma\Sigma^+$、$\Sigma^+\Sigma$ 都是对称对角投影。利用 $U,V$ 正交，四个 Penrose 方程立即成立。

为证唯一性，设 $G$ 也满足四个方程。$AG$ 是对称幂等矩阵，且它在 $\operatorname{Im}A$ 上为恒等映射，故它是到 $\operatorname{Im}A$ 的唯一正交投影，从而
$AG=AA^+$。同理 $GA=A^+A$。因此
\[
G=GAG=GAA^+=A^+AA^+=A^+.
\]
\end{proof}

\begin{corollary}[满秩情形]
若 $A$ 列满秩，则
\[
A^+=(A^{\mathsf T}A)^{-1}A^{\mathsf T}.
\]
若 $A$ 行满秩，则
\[
A^+=A^{\mathsf T}(AA^{\mathsf T})^{-1}.
\]
若 $A$ 是可逆方阵，则 $A^+=A^{-1}$。
\end{corollary}

\begin{proposition}[伪逆的基本性质]
\[
(A^+)^+=A,
\qquad
(A^{\mathsf T})^+=(A^+)^{\mathsf T},
\qquad
\rank(A^+)=\rank(A),
\]
\[
(A^{\mathsf T}A)^+=A^+(A^+)^{\mathsf T},
\qquad
A^+=(A^{\mathsf T}A)^+A^{\mathsf T}
=A^{\mathsf T}(AA^{\mathsf T})^+.
\]
一般情况下 $(AB)^+\ne B^+A^+$；只有在额外秩或子空间条件下才能使用逆序律。
\end{proposition}

\section{伪逆的四个投影与基本子空间}

\begin{theorem}[伪逆投影]
对 $A\in\R^{m\times n}$，
\[
AA^+=P_{\operatorname{Im}A},
\qquad
A^+A=P_{\operatorname{Im}A^{\mathsf T}},
\]
分别是到列空间和行空间的正交投影。同时
\[
I-AA^+=P_{\ker A^{\mathsf T}},
\qquad
I-A^+A=P_{\ker A}.
\]
\end{theorem}
\begin{proof}
由 Penrose 方程，
\[
(AA^+)^2=AA^+AA^+=AA^+,
\qquad
(AA^+)^{\mathsf T}=AA^+.
\]
所以 $AA^+$ 是正交投影。又
$\operatorname{Im}(AA^+)\subseteq\operatorname{Im}A$，而对任意 $y=Ax$，
\[
AA^+y=AA^+Ax=Ax=y,
\]
故其像空间正是 $\operatorname{Im}A$。其余结论同理，或利用正交补关系
$(\operatorname{Im}A)^\perp=\ker A^{\mathsf T}$。
\end{proof}

\begin{pitfall}[幂等投影与正交投影]
$P^2=P$ 只说明 $P$ 是沿某个补空间的投影，这种投影可能是斜投影。只有再满足 $P^{\mathsf T}=P$ 时，$P$ 才是正交投影。给定像空间后，一般投影不唯一，但正交投影唯一。
\end{pitfall}

\begin{proposition}[投影矩阵的谱、秩与迹]
若 $P^2=P$，则：
\begin{enumerate}[label=(\arabic*)]
\item $P$ 的特征值只能是 $0$ 或 $1$；
\item $P$ 可对角化；
\item $\rank P=\tr P$；
\item $\det P\in\{0,1\}$，且 $\det P=1$ 当且仅当 $P=I$。
\end{enumerate}
若 $P\ne0$ 还是正交投影，则 $\lVert P\rVert_2=1$。
\end{proposition}
\begin{proof}
若 $Pv=\lambda v$ 且 $v\ne0$，则
\[
\lambda^2v=P^2v=Pv=\lambda v,
\]
故 $\lambda\in\{0,1\}$。最小多项式整除 $t(t-1)$，其无重根，故 $P$ 可对角化。对角形中 $1$ 的个数同时等于秩和迹，其余结论随即得到。对非零正交投影，其奇异值也只能是 $0,1$，且至少有一个 $1$，故谱范数为 $1$。
\end{proof}

\begin{remark}
非零斜投影仍有谱半径 $1$，但其谱范数可以大于 $1$。这表明特征值不能衡量非正规投影在欧氏范数下的放大效应。
\end{remark}

\section{最小二乘解、最小范数解与解集结构}

\begin{theorem}[全部最小二乘解]
对任意 $A\in\R^{m\times n}$、$b\in\R^m$，最小化问题
\[
\min_x\lVert Ax-b\rVert_2
\]
的全部解为
\[
\boxed{x=A^+b+(I-A^+A)z,
\qquad z\in\R^n.}
\]
所有最小二乘解给出同一预测向量
\[
Ax=AA^+b=P_{\operatorname{Im}A}b.
\]
其中唯一范数最小的解是
\[
x_*=A^+b.
\]
\end{theorem}
\begin{proof}
最小二乘的几何条件是残差与列空间正交，即
\[
Ax=P_{\operatorname{Im}A}b=AA^+b.
\]
因此
\[
A(x-A^+b)=0,
\]
所以 $x-A^+b\in\ker A$。由
$I-A^+A=P_{\ker A}$，得到全部解的表示。

又 $A^+b\in\operatorname{Im}A^{\mathsf T}=(\ker A)^\perp$，因此
\[
\lVert A^+b+(I-A^+A)z\rVert_2^2
=\lVert A^+b\rVert_2^2
 +\lVert(I-A^+A)z\rVert_2^2.
\]
右端在零空间分量为零时唯一取得最小值，故最小范数解为 $A^+b$。
\end{proof}

\begin{corollary}[相容方程的最小范数解]
若 $Ax=b$ 有解，则 $A^+b$ 是全部精确解中范数最小的唯一解。
\end{corollary}

\begin{example}[title={秩不足最小二乘}]
设
\[
A=\begin{pmatrix}1&1\\2&2\\0&0\end{pmatrix},
\qquad
b=\begin{pmatrix}1\\0\\1\end{pmatrix}.
\]
求全部最小二乘解及最小范数解。
\begin{solution}
写成 $A=uv^{\mathsf T}$，其中
\[
u=(1,2,0)^{\mathsf T},
\qquad
v=(1,1)^{\mathsf T}.
\]
秩一矩阵的伪逆为
\[
A^+=\frac{vu^{\mathsf T}}{\lVert u\rVert^2\lVert v\rVert^2}
=\frac1{10}\begin{pmatrix}1&2&0\\1&2&0\end{pmatrix}.
\]
因此
\[
A^+b=\frac1{10}(1,1)^{\mathsf T}.
\]
又 $\ker A=\operatorname{span}\{(1,-1)^{\mathsf T}\}$，故全部最小二乘解为
\[
x=\begin{pmatrix}1/10\\1/10\end{pmatrix}
+t\begin{pmatrix}1\\-1\end{pmatrix},
\qquad t\in\R.
\]
最小范数解对应 $t=0$。
\end{solution}
\end{example}

\section{正则化与伪逆的极限表示}

\begin{theorem}[Tikhonov 正则化]
对 $\lambda>0$，问题
\[
\min_x\left(\lVert Ax-b\rVert_2^2
+\lambda\lVert x\rVert_2^2\right)
\]
有唯一解
\[
x_\lambda=(A^{\mathsf T}A+\lambda I)^{-1}A^{\mathsf T}b.
\]
并且
\[
\lim_{\lambda\downarrow0}x_\lambda=A^+b.
\]
\end{theorem}
\begin{proof}
目标函数的 Hessian 为
$2(A^{\mathsf T}A+\lambda I)$，其正定，故驻点唯一且是全局最小点。对 $A=U\Sigma V^{\mathsf T}$，
\[
x_\lambda
=V\diag\!\left(
\frac{\sigma_i}{\sigma_i^2+\lambda}
\right)U^{\mathsf T}b.
\]
当 $\lambda\downarrow0$ 时，非零奇异值对应系数趋于 $1/\sigma_i$，零奇异值对应系数始终为零，因此极限为 $A^+b$。
\end{proof}

\section{伪逆的稳定性、条件数与逆序律}

\begin{proposition}[伪逆范数]
若 $A$ 的最小非零奇异值为 $\sigma_r(A)$，则
\[
\lVert A^+\rVert_2=\frac1{\sigma_r(A)}.
\]
因此对数据扰动 $b\mapsto b+\delta b$，最小范数解的变化满足
\[
\lVert A^+\delta b\rVert_2
\le\frac{\lVert\delta b\rVert_2}{\sigma_r(A)}.
\]
\end{proposition}
\begin{proof}
$A^+$ 的非零奇异值是
$1/\sigma_1(A),\ldots,1/\sigma_r(A)$。谱范数是最大奇异值，故等于 $1/\sigma_r(A)$。第二式由谱范数的定义立即得到。
\end{proof}

\begin{theorem}[伪逆的连续性边界]
在秩固定的矩阵集合上，映射 $A\mapsto A^+$ 是连续的。但当秩发生变化时，伪逆一般不连续。
\end{theorem}

例如
\[
A_\varepsilon=\begin{pmatrix}1&0\\0&\varepsilon\end{pmatrix}
\longrightarrow
A_0=\begin{pmatrix}1&0\\0&0\end{pmatrix},
\]
但对 $\varepsilon\ne0$，
\[
A_\varepsilon^+
=\begin{pmatrix}1&0\\0&1/\varepsilon\end{pmatrix},
\]
其范数趋于无穷，而
$A_0^+=\diag(1,0)$。该例也说明，接近零的奇异值会使伪逆极度敏感。

\begin{proposition}[伪逆的逆序律]
一般地，
\[
(AB)^+\ne B^+A^+.
\]
若 $A\in\R^{m\times n}$ 列满秩，而
$B\in\R^{n\times p}$ 行满秩，则
\[
(AB)^+=B^+A^+.
\]
\end{proposition}
\begin{proof}
在满秩假设下，
\[
A^+A=I_n,
\qquad
BB^+=I_n.
\]
令 $G=B^+A^+$。则
\[
(AB)G(AB)=A(BB^+)(A^+A)B=AB,
\]
\[
G(AB)G=B^+(A^+A)(BB^+)A^+=G.
\]
又
\[
(AB)G=AA^+,
\qquad
G(AB)=B^+B
\]
均为对称正交投影，故 $G$ 满足四个 Penrose 方程。

反例可取
\[
A=\begin{pmatrix}1&1\end{pmatrix},
\qquad
B=\begin{pmatrix}1\\0\end{pmatrix}.
\]
此时 $AB=(1)$，故 $(AB)^+=(1)$，而
\[
B^+A^+
=\begin{pmatrix}1&0\end{pmatrix}
\frac12\begin{pmatrix}1\\1\end{pmatrix}
=\frac12.
\]
\end{proof}

\section{秩分解与伪逆的紧致构造}

\begin{definition}[秩分解]
若 $A\in\R^{m\times n}$ 的秩为 $r$，且
\[
A=FG,
\]
其中 $F\in\R^{m\times r}$ 列满秩，
$G\in\R^{r\times n}$ 行满秩，则称 $A=FG$ 为 $A$ 的秩分解。
\end{definition}

可从 $A$ 中选取 $r$ 个线性无关列作为 $F$ 的列。由于 $F$ 的列张成 $A$ 的列空间，$A$ 的每一列都可用 $F$ 的列唯一线性表示；把这些坐标依次排成列，即得 $G$。

\begin{theorem}[秩分解给出的伪逆公式]
若 $A=FG$ 是秩分解，则
\[
\boxed{
A^+=G^{\mathsf T}(GG^{\mathsf T})^{-1}
(F^{\mathsf T}F)^{-1}F^{\mathsf T}.}
\]
\end{theorem}
\begin{proof}
由 $F$ 列满秩、$G$ 行满秩，
\[
F^+=(F^{\mathsf T}F)^{-1}F^{\mathsf T},
\qquad
G^+=G^{\mathsf T}(GG^{\mathsf T})^{-1}.
\]
令 $H=G^+F^+$。利用
$F^+F=I_r$、$GG^+=I_r$，
\[
AHA=FGG^+F^+FG=FG=A,
\]
\[
HAH=G^+F^+FGG^+F^+=G^+F^+=H.
\]
又
\[
AH=FF^+,
\qquad
HA=G^+G,
\]
二者分别是到 $\operatorname{Im}F$和
$\operatorname{Im}G^{\mathsf T}$ 的正交投影，因而对称。故 $H$ 满足四个 Penrose 方程，即 $H=A^+$。
\end{proof}

\begin{remark}
秩分解只使用 $r\times r$ 矩阵
$F^{\mathsf T}F$ 和 $GG^{\mathsf T}$ 的逆，因而很适合手算低秩矩阵的伪逆。数值计算中仍应避免在病态问题上直接形成 Gram 矩阵，此时 SVD 更稳定。
\end{remark}

\section{参数线性方程的微分}

\begin{theorem}[隐式线性方程的微分]
设 $A=A(\theta)$ 在所考察参数附近可逆，
$b=b(\theta)$，且 $x=x(\theta)$ 由
\[
A(\theta)x(\theta)=b(\theta)
\]
唯一确定。则
\[
\boxed{dx=A^{-1}(db-(dA)x).}
\]
若 $b$ 为常向量，则
\[
dx=-A^{-1}(dA)x.
\]
\end{theorem}
\begin{proof}
对 $Ax=b$ 求微分，由乘积法则得
\[
(dA)x+A\,dx=db.
\]
移项后左乘 $A^{-1}$，得
\[
dx=A^{-1}db-A^{-1}(dA)x.
\]
\end{proof}

\begin{corollary}[具体参数的导数]
若 $A(t)x(t)=b(t)$，则
\[
x'(t)=A(t)^{-1}\bigl(b'(t)-A'(t)x(t)\bigr).
\]
该公式不需要先把 $x(t)=A(t)^{-1}b(t)$ 展开后再对逆矩阵求导。
\end{corollary}

\begin{pitfall}[可逆性是公式的关键条件]
若 $A(\theta)$ 在某点奇异，方程的解可能不存在、不唯一或对参数不连续，上述公式不能直接使用。使用伪逆进行求导还需要秩在邻域内保持不变，属于更高级的矩阵扰动问题。
\end{pitfall}

\section{一般广义逆的全部参数化}

\begin{theorem}[所有内逆的参数化]
固定 $A$ 的一个反身广义逆 $R$，即
\[
ARA=A,
\qquad
RAR=R.
\]
例如可取 $R=A^+$。则所有满足 $AGA=A$ 的矩阵 $G$ 恰好可表为
\[
\boxed{G=R+Y-RAYAR,}
\]
其中 $Y\in\R^{n\times m}$ 任意。
\end{theorem}
\begin{proof}
对任意 $Y$，
\begin{align*}
AGA
&=ARA+AYA-ARAYARA\\
&=A+AYA-AYA=A,
\end{align*}
其中使用了 $ARA=A$。因此公式中的每个 $G$ 都是广义逆。

反之，设 $G$ 满足 $AGA=A$。在参数公式中取 $Y=G$，则
\[
R+G-RAGAR
=R+G-R(AGA)R
=R+G-RAR=G.
\]
故任意广义逆都被该参数化包含。
\end{proof}

\begin{remark}
参数矩阵 $Y$ 不唯一；该公式的作用是刻画全部广义逆的集合，而不是给每个 $G$ 分配唯一坐标。若基准矩阵只满足 $ARA=A$ 而不满足 $RAR=R$，上述反向覆盖证明不再成立。
\end{remark}

\begin{method}[矩阵求导与伪逆题的方法选择]
标量对矩阵求导时，优先用 $df=\tr((\nabla_Xf)^{\mathsf T}dX)$；复合映射则先写每一层的微分。参数方程 $Ax=b$ 可直接对等式求微分，避免展开逆矩阵。求伪逆时，一般矩阵用 SVD，低秩手算可用秩分解，满列秩或满行秩时再使用简化公式。求最小二乘解集时，先写 $A^+b$，再加上 $\ker A$ 分量；要求最小范数时，去掉零空间分量。
\end{method}

\begin{summary}[本章知识链]
全微分把非线性变化的一阶部分表示为线性映射；迹和 Frobenius 内积把标量对矩阵的导数统一成矩阵梯度。Vec 与 Kronecker 积则把矩阵映射还原为普通线性映射，特征值扰动公式把谱量的变化与特征投影相联系，参数方程微分则把逆矩阵扰动化为线性方程。加权和约束最小二乘将正规方程推广到非标准内积与可行集上。一般广义逆解决非方阵和奇异矩阵的代数反解问题，秩分解和全部参数化分别给出紧致构造与不唯一性的完整描述；Moore--Penrose 逆则通过对称性选出唯一的正交几何解。最小非零奇异值同时控制伪逆范数、扰动放大和秩变化附近的不稳定性。
\end{summary}
